% ---------------------------------------------------------------------
\section{$B\to D\pi$：领头幂次的全阶因子化}\label{sec:heavy-bdpi}

在“I 类”衰变 $\bar B^0\to D^+\pi^-$ 中，旁观反夸克进入 $D^+$，$\pi^-$ 由弱顶点直接产生（\cref{fig:heavy-bdpi}）。按色透明性（colour transparency）\cite{Bjorken:1988kk}，快速的 $\pi$ 形成之初是小的色单态偶极子，软胶子看不到其色荷，振幅近似朴素因子化\cite{Dugan:1990de,Politzer:1991au}；BBNS 作了两圈分析\cite{BenekeEtAl2000}，BPS 给出 SCET 全阶论证\cite{BauerPirjolStewart2001}。

%%INSERT fig_bdpi.tex

\paragraph{朴素因子化与色透明性。}
$\mathcal H_W=\frac{4G_F}{\sqrt2}V_{cb}V_{ud}^*[C_{\mathbf 0}O_{\mathbf 0}+C_{\mathbf 8}O_{\mathbf 8}]$，其中\cite[式~(2)]{BauerPirjolStewart2001}
\begin{equation}\label{eq:heavy-O08}
O_{\mathbf 0}=\big[\bar c\gamma^\mu P_Lb\big]\big[\bar d\gamma_\mu P_Lu\big],\qquad
O_{\mathbf 8}=\big[\bar c\gamma^\mu P_Lt^ab\big]\big[\bar d\gamma_\mu P_Lt^au\big],\qquad P_L=\tfrac12(1-\gamma_5).
\end{equation}
% VF-BQ-6
在 $\bar B$ 静止系中取无质量的 $\pi$ 沿 $+z$ 运动，$p_\pi=\sqrt2E_\pi n_A$，$E_\pi=\frac{m_B^2-m_D^2}{2m_B}\approx2.31$\,GeV（$m_B\approx5.28$\,GeV，$m_D\approx1.87$\,GeV\cite{ParticleDataGroup:2024cfk}）。真空插入近似给出 $\langle\pi^-|\bar d\gamma_\mu P_Lu|0\rangle\langle D^+|\bar c\gamma^\mu P_Lb|\bar B^0\rangle$：由\cref{def:excl-da}的 $\langle0|\bar d\gamma^\mu\gamma_5u|\pi^+(p)\rangle=\ii f_\pi p^\mu$（共轭加同位旋）与宇称，前一因子为 $\frac{\ii}{2}f_\pi p_{\pi\mu}$；后一因子与 $q=p_\pi$ 收缩时 $F_+$ 项消去，剩 $\frac12(m_B^2-m_D^2)F_0(0)$，$F_0(0)=F_+(0)\equiv F^{B\to D}(0)$。于是
\begin{equation}\label{eq:heavy-naive}
\langle D^+\pi^-|O_{\mathbf 0}|\bar B^0\rangle\big|_{\rm naive}=\frac{\ii}{2}f_\pi\cdot\frac12(m_B^2-m_D^2)F^{B\to D}(0)=\frac{\ii}{2}\,m_BE_\pi f_\pi\,F^{B\to D}(0),
\end{equation}
即下面 $T=1$ 时的归一化 $N$（用 $(V-A)\otimes(V-A)$ 算符时大 $4$ 倍）；$O_{\mathbf 8}$ 的真空插入 $\propto\Tr t^a=0$。软胶子（$k\sim\Lqcd$）与出射夸克 $p_1=xp_\pi+l$、反夸克 $p_2=\bar xp_\pi-l$（$\bar x=1-x$，横向 $|l|\sim\Lqcd$）的程函耦合，对色单态对给出同一色矩阵而符号相反：
\begin{equation}\label{eq:heavy-dipole}
J^\mu\propto t^a\Big[\frac{p_1^\mu}{p_1\cdot k}-\frac{p_2^\mu}{p_2\cdot k}\Big]
=t^a\Big(\frac1x+\frac1{\bar x}\Big)\frac{1}{p_\pi\cdot k}\Big[l^\mu-p_\pi^\mu\,\frac{l\cdot k}{p_\pi\cdot k}\Big]+\order{l^2}.
\end{equation}
$l=0$ 时 $J^\mu\equiv0$（程函因子只依赖方向）；余项是守恒的色偶极子（$k_\mu J^\mu=0$），相对单条线为 $\order{\Lqcd/(x\bar xm_b)}$。所以软胶子在领头幂次退耦——\emph{端点 $x\bar x\sim\Lqcd/m_b$ 除外}，这正是需要端点压低假设的原因（脚本 \texttt{s1\_\allowbreak bdpi.py}）。

\begin{theorem}[$B\to D\pi$ 的全阶因子化 \tagA]\label{thm:heavy-bdpi}
设 $Q\in\{m_b,m_c,E_\pi\}$ 都远大于 $\Lqcd$，$\lambda=\Lqcd/Q$，并设（端点压低假设）介子光锥波函数的端点组态——$\pi$ 中一个价夸克带走几乎全部能量的组态，以及 $B$、$D$ 中旁观者动量 $\gg\Lqcd$ 的组态——被 $\lambda^{a}$（某个 $a>0$）压低。则对 $\bar B^0\to D^+\pi^-$ 与 $B^-\to D^0\pi^-$，在 $\lambda$ 的领头幂次、对 $\as$ 的所有阶，
\begin{equation}\label{eq:heavy-bdpi}
\langle D\pi|O|B\rangle=N\,F^{B\to D}(0)\int_0^1\dd x\;T_O(x,\mu)\,\phi_\pi(x,\mu)\,\big[1+\order{\lambda}\big],\qquad N=\frac{\ii}{2}\,m_BE_\pi f_\pi ,
\end{equation}
% VF-BQ-6
其中 $O=O_{\mathbf 0}$ 或 $O_{\mathbf 8}$（\cref{eq:heavy-O08}），$F^{B\to D}(0)$ 是 $q^2=0$ 处的 $B\to D$ 形状因子，$\phi_\pi$ 是\cref{def:excl-da}的分布振幅（$x$ 为 $\pi^-$ 中 $d$ 夸克的动量分数；由 G 宇称 $\phi_\pi(x)=\phi_\pi(1-x)$），$T_O$ 是可微扰计算的短程系数，$T_{O_{\mathbf 0}}=1+\order{\as}$，$T_{O_{\mathbf 8}}=\order{\as}$。取 $F^{B\to D}(0)$ 还是 $F^{B\to D}(m_\pi^2)$ 只差幂次压低的项。
非微扰地位：\cref{asm:intro-bridge}（MA1、MA2）。
% errtype: E-S
\uses{def:excl-da, thm:scet-decoupling, asm:intro-bridge}
\end{theorem}

\begin{proofsketch}
按\cite{BauerPirjolStewart2001}，在 SCET 与 HQET 中进行。(1) 剩下的模式是 HQET 重夸克与软模式（$\sim Q\lambda$）、$\pi$ 的共线夸克 $\xi$（$\sla n_A\xi=0$）与胶子（$Q(1,\lambda^2,\lambda)$）、超软胶子；它们穷尽领头区域是 SCET 的第一前提（E-S 型陈述）。(2) 共线胶子接到重夸克上使其离壳，积出后求和为沿 $n_B$ 的共线 Wilson 线 $W$：$O_{\mathbf 0,\mathbf 8}$ 匹配到 $Q_{\mathbf 0}=[\bar h_{v'}\Gamma_hh_v][\bar\xi\Gamma_\ell\xi]$、$Q_{\mathbf 8}=[\bar h_{v'}\Gamma_h(W^\dagger t^aW)h_v][\bar\xi\Gamma_\ell t^a\xi]$（对色单态 $W^\dagger W=1$）。(3) 超软胶子与重夸克的耦合被 $\lambda$ 压低，与共线场的耦合由 BPS 场重定义消去（\cref{thm:scet-decoupling}，[R, R$_{\rm def}$]）；$\order{\lambda^0}$ 只有 $n_A\cdot A_s$ 与共线粒子耦合，且软动量只以 $n_A\cdot k$ 进入共线传播子，Ward 恒等式把软胶子提出为沿 $n_A$ 的软 Wilson 线 $S$：对 $Q_{\mathbf 0}$，$S^\dagger S=1$（\cref{eq:heavy-dipole}的全阶形式）；对 $Q_{\mathbf 8}$，$\Tr[MW^\dagger t^aW]=M\Tr t^a=0$\cite[式~(7)]{BauerPirjolStewart2001}。(4) 共线规范不变性把硬修正限制为 $[\bar h_{v'}\Gamma_hh_v][\bar\xi W\,C(\bar{\mathcal P}_+)\Gamma_\ell W^\dagger\xi]$，$C$ 只依赖于夸克与反夸克大光锥分量之差；对 $\pi$ 态 $\bar{\mathcal P}_+\to(2x-1)p_\pi^+$，矩阵元是 $\phi_\pi$ 的 Fourier 变换（\cref{def:excl-da}），故 $T(x,\mu)=C(\mu,(2x-1)p_\pi^+)$。(5) 端点压低使弱湮灭、硬旁观者及旁观者进入 $\pi$ 的拓扑幂次压低（如硬旁观者图 $\sim G_F\as m_b\Lqcd^2$，领头项 $\sim G_Fm_b^2\Lqcd$\cite{BenekeEtAl2000}）。证明不需要 $\pi$ 由 $q\bar q$ Fock 态主导；两圈核见\cite{Huber:2016xod}。
\end{proofsketch}

\paragraph{地位与缺口}
全阶论证（[A]）。输入之一是 SCET 的前提（\cref{rem:scet-assumptions}）：模式穷尽领头区域（即\cref{ch:ir}的区域分析，强子层次上为\cref{asm:intro-bridge}(Br1)）以及 Glauber 模式与零箱减除的处理，只给出 E-S 型陈述。输入之二是端点压低假设，它是关于强子波函数的非微扰陈述，BPS 明确列为假设\cite{BauerPirjolStewart2001}，没有证明。BBNS 在两圈显式证明了红外奇异性的相消，并写明“the following does not accomplish an all-order ‘proof’”\cite[第~5~节]{BenekeEtAl2000}。没有已知的领头幂次反例。升级途径：对本过程证明模式完备性，并证明端点压低。
